\begin{lemma}
Let $1<p<\infty$ and let $\kappa(p)$ be the right-hand side of equation (9),
as in \noderef{KAlphaExp}.  Then there is exactly one pair $(k,\alpha)$
consisting of an integer $k\ge1$ and a real number $\alpha$ with
$0<\alpha\le1$ such that
\[
  k+\alpha=\kappa(p).
\]
\end{lemma}
\begin{proof}
Fix $p>1$ and write $v=\kappa(p)$ and $q=\dfrac1{p-1}$, so that $q>0$ and, by
\noderef{KAlphaExp},
\[
  6v=7+q+\sqrt{1+14q+q^2}.
\]

\emph{Step 1 ($v>\tfrac43$, in particular $v>1$).}  Since $q>0$ we have
$14q>0$ and $q^2\ge0$, so $1+14q+q^2>1$.  The square root is strictly
increasing on $[0,\infty)$ and $\sqrt1=1$, so $\sqrt{1+14q+q^2}>1$.  Hence
\[
  6v=7+q+\sqrt{1+14q+q^2}>7+0+1=8 ,
\]
so $v>\tfrac86=\tfrac43>1$.

\emph{Step 2 (existence).}  Let $m=\lceil v\rceil$ be the least integer
$\ge v$, and set
\[
  k=m-1,\qquad \alpha=v-k .
\]
Since $v>1$, the integer $1$ is not $\ge v$, so $m\ge2$ and therefore
$k=m-1\ge1$ is an integer $\ge1$.  By definition of the ceiling we have
$v\le m$, and also $m-1<v$, for otherwise $m-1$ would be an integer $\ge v$
smaller than $m$.  Thus
\[
  k<v\le k+1 ,
\]
and hence $\alpha=v-k$ satisfies $0<\alpha\le1$.  Finally
$k+\alpha=k+(v-k)=v$.  So the pair $(k,\alpha)$ has all the required
properties.

\emph{Step 3 (uniqueness).}  Suppose $(k,\alpha)$ and $(k',\alpha')$ both
consist of an integer $\ge1$ and a real number in $(0,1]$, and that
\[
  k+\alpha=v=k'+\alpha' .
\]
Then $k-k'=\alpha'-\alpha$.  From $0<\alpha\le1$ and $0<\alpha'\le1$ we get
\[
  -1\le-\alpha<\alpha'-\alpha<\alpha'\le1,
\]
so $\alpha'-\alpha$ lies in the open interval $(-1,1)$.  Hence the integer
$k-k'$ lies in $(-1,1)$, which forces $k-k'=0$, i.e.\ $k=k'$.  Then
$\alpha=v-k=v-k'=\alpha'$.  So the two pairs coincide.
\end{proof}
